\section{Confronting theory with robotic experiments}
\label{sec:theory}

Robotic active matter offers an unusual opportunity to test
theories of collective motion. In a conventional experiment,
the microscopic interactions are often determined by material
properties and cannot be varied independently. Robots permit
a different approach. We can prescribe alignment, alter
feedback, modify interaction networks, or change mechanical
coupling while observing the resulting collective dynamics.

Yet, this advantage introduces an important ambiguity.
A theory implemented directly as a control algorithm is
not necessarily a theory independently validated by
experiment. The hardware may reproduce the intended
interaction while the collective observables remain
untested. Conversely, a model may describe an experimental
pattern accurately without identifying the mechanism
that generated it.

We therefore distinguish implementation, qualitative
consistency, quantitative agreement, and predictive
validation. An implementation demonstrates that a
specified microscopic rule can operate in a physical
system. Qualitative agreement indicates that the
model reproduces a relevant pattern or trend.
Quantitative agreement requires comparison of numerical
observables. A predictive test is more demanding:
parameters must be constrained independently, and
the model must successfully describe outcomes
not used to construct or calibrate it.

The evidence base contains 491 studies in
the October 17, 2026 snapshot. The corrected
extraction table contains 491 records, and
the framework counts reported below are
calculated from this evidence base.
The citation keys remain those supplied
in the extraction table.

Theoretical frameworks were coded using a multi-label
vocabulary. A study could compare several models,
including a continuum theory and an agent-based
simulation. Consequently, the framework totals
cannot be added to obtain a unique number of
theory-tested studies.

Figure~\ref{fig:theory} summarizes these comparisons.
The figure distinguishes the extracted frameworks
and their reported levels of theory--experiment
agreement. It describes what was compared, not
whether a model was tested predictively.

\begin{figure}[t]
\centering
\includegraphics[width=\columnwidth]{figures/subtaxonomy_theory}
\caption{Theoretical frameworks compared with robotic
experiments in the extracted evidence base.
Framework categories are multi-label and therefore
not mutually exclusive. Bars are stacked by the
reported type of agreement: quantitative,
qualitative, or none. An agreement label does
not establish independent parameter calibration
or predictive validation. The figure represents
the original extraction snapshot; the corrected
491-study corpus differs by one merged microscopic
record whose identity is not specified here.}
\label{fig:theory}
\end{figure}

The distinction between agreement and predictive
validation guides the following analysis. For each
framework, we ask which assumptions can be realized
experimentally, which observables were compared,
how much quantitative evidence is available,
and where the implementation itself limits the
strength of the theoretical conclusion.

\subsection{Vicsek-type models and alignment theories}
\label{sec:theory-vicsek}

Vicsek-type models provide a natural starting point
because their elementary ingredients resemble the
capabilities of programmable robots. Self-propelled
agents move while adjusting their orientations
according to neighboring agents and fluctuations.
The collective dynamics can produce polarization
and related forms of coordinated motion
\cite{Vicsek1995novel}.

In a representative discrete-time formulation,
\begin{equation}
\theta_i(t+\Delta t)=
\operatorname{Arg}
\left[
\sum_{j\in\mathcal N_i}
e^{\mathrm{i}\theta_j(t)}
\right]
+\eta\xi_i(t),
\label{eq:theory-vicsek}
\end{equation}
where $\mathcal N_i$ denotes the interaction
neighborhood and $\eta\xi_i$ describes angular
noise. The position update follows from
self-propulsion along the selected direction.

This equation captures a simple idea:
local alignment can generate global order.
Yet, its physical consequences are not limited
to the appearance of aligned trajectories.
The theory makes statements about fluctuations,
density-dependent organization, collective modes,
and the behavior of increasingly large systems
\cite{Toner1995longrange}
\cite{Vicsek1995novel}.

A robot swarm can implement an alignment rule
with considerable fidelity. However, the
implemented rule is generally accompanied by
sensor noise, communication delays, finite
turning rates, collision avoidance, and
environmental constraints. These effects
can change the resulting dynamics.

Hence, the central experiment is not simply
whether robots align. It is whether their
measured collective response agrees with
the predictions of a specified model once
the actual implementation is characterized.

\paragraph*{Experimental realizations and observables.}

The extraction identifies 20 studies comparing
with a Vicsek-type framework. Ten involve
wheeled or legged swarm robots, five belong
to the residual platform category, three
use feedback-controlled colloids, and one
uses vibration-driven robots. One additional
record has no platform classification.

The experiments cover programmed alignment,
collective motion, orientation-dependent
interactions, and related collective states
\cite{Turgut2008selforganizedflocka,Tarcai2011patternstransition,Vsrhelyi2014outdoorflockingand,Balzs2020adaptiveleadership,Fernando2021onlineflockingcont,Wang2022optiswarmopticals,Xiao2023perceptionofmotion,Zheng2024bodyorientationcha,Jianchao2024emergentchiralitya,Verdoucq2025flockingphasetrans,Zheng2026nestedinteractionn,Dizenhaus2026anonymousleadership,Zhang2026noiseresilientimpl,Yasuda2014selforganizedflockb}.

Their extracted observables include quantities
related to group orientation, polarization,
collective trajectories, spatial organization,
and changes in collective motion. These
measurements are appropriate for evaluating
alignment models, although they constrain
different aspects of the theory.

Polarization, for example, characterizes
global directional coherence. It cannot
by itself distinguish a homogeneous
ordered state from spatially heterogeneous
collective organization. Spatial correlations
and density-dependent measurements provide
additional information.

Similarly, a trajectory showing successful
collective turning establishes that
the implemented interaction produces
a coordinated response. It does not
necessarily determine whether fluctuations
or relaxation obey the proposed model.

The feedback-controlled colloidal experiments
provide another experimental route
\cite{Chen2024synchronousandfull,Chen2022collectiveresponse,Chen2024persistentandrespo}.
Their importance is that effective
orientation-dependent interactions can
be investigated at the particle scale.
Such systems may provide more direct
access to collective order parameters,
but the implemented feedback law
must still be separated from
the resulting physical predictions.

Other comparisons investigate collective
motion and ordering through nonstandard
or asymmetric interactions
\cite{Fruchart2021nonreciprocalphase,Giomi2013swarmingswirlinga,Liu2025collectivemotionmo}.
These studies broaden the comparison
beyond a single implementation of
local alignment.

However, departures from the basic
Vicsek assumptions should not be
treated merely as experimental noise.
They may define a different model.
For example, nonreciprocal interactions,
delayed information, and persistent
turning dynamics can change the
appropriate theoretical description.

\paragraph*{Quantitative agreement and its limits.}

Only four of the 20 Vicsek-type studies
have an extracted quantitative agreement
classification
\cite{Chen2024synchronousandfull,Xiao2023perceptionofmotion,Dizenhaus2026anonymousleadership,Chen2024persistentandrespo}.
Fourteen are classified as qualitative
comparisons, while two have no
agreement comparison recorded.

The risk-of-bias assessment reinforces
this distinction. Four studies receive
low-risk ratings in D6, eleven have
some concerns, and five are high risk.
The low-risk subset is therefore
considerably smaller than the complete
literature implementing or discussing
Vicsek-type dynamics.

For replication, eight studies
receive low-risk D1 ratings and
twelve have some concerns.
Independent-run counts are recorded
for eleven studies. Explicit code
availability occurs in two and
data availability in three.
All 20 studies were assessed
on a full-text basis.

These results establish that Vicsek-type
models have been realized and compared
with robotic experiments in several
physical implementations. Yet, most
reported comparisons remain qualitative.

This is particularly consequential
because the alignment rule itself
is easy to implement. The experiment
can reproduce the mathematical
update mechanism while failing
to test the collective predictions
that make the theory physically
interesting.

A stronger validation would characterize
the implemented alignment interaction,
measure the dependence of polarization
and correlations on relevant parameters,
and compare the resulting behavior
with predictions obtained before
observing the experimental outcome.

Finite-size analysis is especially
important. A small group of aligned
robots is compatible with many
possible interaction mechanisms.
Determining whether the collective
obeys a particular ordering theory
requires measurements that distinguish
between competing models.

The evidence therefore supports
the feasibility of Vicsek-type
robotic realizations more strongly
than it supports universal conclusions
about Vicsek-type phase behavior.
The principal opportunity is to
use programmable alignment to
measure collective fluctuations,
responses, and scaling laws
under controlled physical conditions.

\subsection{Active Brownian particle models}
\label{sec:theory-abp}

Active Brownian particle models provide
a different theoretical description.
Instead of prescribing explicit
alignment, they treat individual
agents as persistently self-propelled
particles whose orientations fluctuate.
Interactions may then arise through
steric contact, external potentials,
or additional physical coupling
\cite{Bechinger2016active}
\cite{Bechinger2016active}.

In a simple two-dimensional formulation,
\begin{align}
\dot{\mathbf r}_i &=
v_0\mathbf e(\theta_i)
+\mu\mathbf F_i
+\sqrt{2D_t}\,\boldsymbol{\xi}_i(t),
\\
\dot{\theta}_i &=
\sqrt{2D_r}\,\eta_i(t),
\label{eq:theory-abp}
\end{align}
where $v_0$ is the self-propulsion speed,
$D_t$ and $D_r$ are translational and
rotational diffusivities, and
$\mathbf F_i$ denotes interactions
or external forces.

The physical interpretation is direct.
Even without an explicitly programmed
alignment rule, persistent propulsion
and interactions can generate
nonequilibrium collective dynamics.

This makes active Brownian particle
models particularly suitable for
mechanically driven robotic systems.
Yet, the assumptions of constant
self-propulsion, simple stochastic
reorientation, and memoryless noise
must be checked experimentally.

Real vibration-driven robots may exhibit
intermittent propulsion, correlated
rotational motion, asymmetric body
geometry, and substrate-dependent
friction. These effects may be
essential to the observed collective
behavior rather than negligible
deviations from the ideal model.

\paragraph*{Experimental implementations.}

The extraction identifies 28 studies
with active Brownian particle comparisons.
Eighteen use vibration-driven robots,
five wheeled or legged robots,
three robotic micro- or nano-agents,
and two feedback-controlled colloids.

The substantial representation of
vibration-driven systems distinguishes
this framework from Vicsek-type models.
These experiments investigate
mechanically generated persistent
motion, aggregation, transport,
confinement, and collective
reorganization
\cite{Welch2018aggregatesuppressio,Balda2022playingwithactive,Fersula2026aggregatingswarmst,Edwards2026robusttopologically,Musacchio2026activeembracemente,Hoek2026pistonlikeinformat,Arbel2025amechanicalroutef,Fersula2026leveragingdesignfo,MuozObreque2026dynamicsofmagnetic,Musacchio2026fluidizationinduced,Callegari2023playingwithactive,Balda2024playingwithactive,Li2021programmingactivec,Siebers2023exploitingcompositi,Son2024emergentfunctional,Zion2023morphologicalcomput,Wagner2026geometryinducedtra,Zion2021distributedonline}.

These systems provide a useful
experimental connection between
microscopic propulsion statistics
and collective observables.
Unlike a programmed alignment rule,
mechanical interactions need not
be prescribed explicitly in
the individual controllers.

The relevant measurements include
individual trajectories,
mean-squared displacement,
rotational dynamics,
density distributions,
clustering statistics,
and confinement-dependent
transport. These quantities
can constrain different
components of the active
Brownian description.

For example, the isolated-agent
mean-squared displacement
can characterize persistence.
Collective measurements
can then reveal how
interactions modify
that baseline behavior.

This separation is fundamental.
A model fitted directly to
collective transport data
may reproduce the observations
without independently establishing
the propulsion statistics.
By contrast, parameters measured
from isolated-agent trajectories
provide a stronger basis
for predicting collective motion.

Wheeled or legged experiments
offer a complementary realization
through programmed or controlled
persistent motion
\cite{Paramanick2025spontaneousemergenc,Barone2024experimentalandnum,Mijalkov2016engineeringsensoria,Mayya2019nonuniformrobotde,Hao2022controllingcollisio}.
Such systems can vary effective
motility and interaction rules
more independently than purely
mechanical particles.

However, they may also introduce
feedback effects not represented
in the standard active Brownian
equations. Thus, the comparison
must consider whether the
robotic dynamics actually
approximate the assumed
stochastic process.

Microscopic and colloidal systems
extend the framework toward
field-mediated and feedback-controlled
interactions
\cite{Rivas2023spatialpatterningo,Raghavendra2026programmablepersist,Desgranges2023microswimmersunder,Knippenberg2022roleofcohesionin,Guo2026afeedbackcontrolle}.
These experiments are relevant
because persistent particle motion
and effective interactions can
be characterized at smaller
length scales.

Yet, common external forcing
and hydrodynamic coupling may
produce correlations absent
from a simple active Brownian
particle model. Quantitative
agreement must therefore be
interpreted relative to the
actual interaction mechanism.

\paragraph*{Quantitative evidence.}

Nineteen of the 28 studies report
quantitative theory--experiment
agreement, while nine report
qualitative agreement. This is
a markedly higher quantitative
fraction than in the Vicsek-type
category.

The quantitative subset includes
mechanically driven systems,
programmable colloids, and
mobile robotic experiments
\cite{Balda2022playingwithactive,Musacchio2026activeembracemente,Arbel2025amechanicalroutef,Knippenberg2022roleofcohesionin,Raghavendra2026programmablepersist,Desgranges2023microswimmersunder,Guo2026afeedbackcontrolle,Paramanick2025spontaneousemergenc,Musacchio2026fluidizationinduced,Barone2024experimentalandnum,Mijalkov2016engineeringsensoria,Hao2022controllingcollisio,Callegari2023playingwithactive,Balda2024playingwithactive,Li2021programmingactivec,Siebers2023exploitingcompositi,Zion2023morphologicalcomput,Wagner2026geometryinducedtra,Zion2021distributedonline}.

For D6, twelve studies are
rated low risk, fourteen
have some concerns, and
two are high risk.
Thus, quantitative
agreement is comparatively
well represented, but
some reported comparisons
retain methodological
limitations.

Replication evidence is
less favorable. Only six
studies receive low-risk
D1 ratings, twenty have
some concerns, and two
are high risk. Independent-run
counts are recorded for
eight studies.

Explicit code availability
occurs in six studies,
and data availability
in seven. One study
was assessed from its
abstract alone. Its
contribution to claims
requiring detailed
experimental parameter
characterization is
therefore less secure.

The quantitative evidence
supports the usefulness
of active Brownian particle
models for describing
several robotic experiments.
Yet, the extraction does
not systematically
distinguish fitted
parameters from
independently measured
parameters.

Consequently, nineteen
quantitative comparisons
must not be interpreted
as nineteen independent
predictive validations
of the active Brownian
framework.

\paragraph*{What would distinguish competing descriptions?}

The most informative tests
begin with isolated-agent
dynamics. An experimental
determination of propulsion
speed, orientation persistence,
and stochastic fluctuations
provides a reference against
which collective behavior
can be predicted.

Subsequently, particle
density and confinement
can be varied while
retaining those measured
parameters. The resulting
collective observables
can be compared with
predictions of the model.

For clustering, this
would help distinguish
motility-induced
organization from
direct attraction.
For transport, it
would establish whether
changes in mobility
follow from persistence
and interactions.
For boundary accumulation,
it would test whether
the observed density
profile follows from
the experimentally
measured single-agent
dynamics.

A particularly demanding
comparison would use
one parameterization
to predict several
observables simultaneously.
A model reproducing
mean-squared displacement,
density correlations,
and confinement response
without separate fitting
for each measurement
would provide stronger
evidence than agreement
with one trajectory
distribution.

Another important
question concerns
the validity of
effective noise models.
If vibration-driven
robots possess
correlated or
non-Gaussian fluctuations,
the standard active
Brownian description
may reproduce
long-time transport
while failing
at shorter times.

This would not
necessarily invalidate
the model.
It would identify
its domain of
applicability.
Such boundaries
are scientifically
valuable because
they reveal when
additional physical
degrees of freedom
must be included.

We therefore regard
active Brownian particle
models as one of
the better-supported
quantitative frameworks
in the robotic
active-matter evidence
base. Their experimental
breadth is narrower
than that of general
agent-based simulation,
but their connection
to measurable
single-particle
dynamics makes
them particularly
suitable for
discriminating tests.

The remaining challenge
is to move from
agreement with
selected observables
toward parameter-constrained
predictions of
collective behavior.

\subsection{Run-and-tumble descriptions}
\label{sec:theory-run-tumble}

Run-and-tumble models describe active motion as alternating
episodes of persistent displacement and reorientation.
They provide a useful alternative to continuous
rotational diffusion. In the simplest formulation,
an agent moves at approximately constant speed
during a run and changes its direction through
discrete reorientation events
\cite{Tailleur2008runtumble}.

A characteristic observable is the distribution
of run durations. If reorientations occur as
a Poisson process with rate $\alpha$,
the probability that a run persists
longer than time $t$ is
\begin{equation}
S(t)=\exp(-\alpha t).
\label{eq:theory-run-tumble}
\end{equation}
This expression is a model prediction,
not a general property of every
run-and-tumble robotic system.
Nonexponential waiting times,
correlated turning angles,
and state-dependent reorientation
require different descriptions.

The distinction from active Brownian
motion concerns the microscopic
statistics of orientation changes.
Both descriptions can generate
persistent trajectories and
enhanced long-time spreading.
Therefore, agreement with
mean-squared displacement alone
may not distinguish them.

We need measurements of turning
events, orientation correlations,
and run-duration distributions.
The physical interpretation
becomes particularly informative
when these quantities can be
measured independently of
collective behavior.

Only three studies in the extraction
explicitly compare with a
run-and-tumble framework.
Two involve vibration-driven
robots and one a robotic
micro- or nano-agent system
\cite{Noirhomme2024brainbotsassmarta,Raghavendra2026programmablepersist,Balda2024playingwithactive}.

This is a small evidence base
relative to both Vicsek-type
and active Brownian comparisons.
It nevertheless provides
an important conceptual link
between robotic control
and stochastic motility.

A robot can realize
discrete reorientation
through programmed
motor commands,
changes in actuation,
or interaction-mediated
turning. The resulting
trajectories can then
be compared with
predictions based
on run and tumble
statistics.

Yet, a programmed
run-and-tumble sequence
is not automatically
a quantitative test
of run-and-tumble
theory. It may simply
implement the assumed
microscopic dynamics.
The stronger experiment
asks whether collective
transport, response,
or spatial organization
follows from those
independently characterized
statistics.

The three studies include
two quantitative comparisons
and one qualitative comparison.
Two receive low-risk
D6 ratings, while one
is rated high risk.

For replication,
all three have
some concerns in D1.
An independent-run
count is recorded
for one study.
One explicitly reports
code availability,
and one reports
data availability.
All three were
assessed on full text.

These results indicate
that run-and-tumble
descriptions have
been tested quantitatively
in a limited set
of robotic systems.
However, the small
number of implementations
does not permit a
strong claim about
cross-platform robustness.

The main opportunity
is experimental
discrimination between
continuous orientational
diffusion and discrete
reorientation.
Robots are particularly
well suited to this
comparison because
different turning
statistics can be
implemented on
otherwise similar
hardware.

A decisive experiment
would first measure
the complete single-agent
orientation process.
It would then examine
whether the same
parameters predict
collective observables
under changes in
density, interaction,
and confinement.

The comparison should
include observables
sensitive to the
distribution of
reorientation events,
rather than only
their average rate.
This would reveal
when run-and-tumble
and active Brownian
models are effectively
equivalent and when
their different
microscopic assumptions
produce experimentally
distinguishable results.

Hence, the current
evidence supports
the feasibility
of robotic run-and-tumble
experiments, but
general conclusions
about collective
behavior remain
premature. The
framework is a
particularly useful
candidate for
controlled comparative
experiments rather
than a broadly
validated explanation
of existing robotic
collectives.

\subsection{Kinetic and hydrodynamic theories}
\label{sec:theory-kinetic}

Kinetic and hydrodynamic theories provide the connection
between microscopic interactions and collective fields.
Instead of tracking every agent separately, we describe
the evolution of quantities such as density, polarization,
velocity, and stress. This is a central theoretical
step in active-matter physics. It permits predictions
about collective transport, fluctuations, stability,
and spatial organization that cannot be obtained
from the motion of one agent alone
\cite{Toner1995longrange}
\cite{Marchetti2013hydrodynamics}
\cite{Marchetti2013hydrodynamics}.

For a conserved population of agents, the density
$\rho(\mathbf r,t)$ obeys a continuity equation,
\begin{equation}
\partial_t\rho+
\nabla\cdot\mathbf J=0,
\label{eq:theory-continuity}
\end{equation}
where $\mathbf J$ is the particle current.
A common active contribution takes the form
$\mathbf J_{\mathrm a}=v_0\rho\mathbf p$,
with $\mathbf p$ denoting local polarization.
Further equations specify the dynamics of
orientation, velocity, or stress.

The difficulty lies in constructing and testing
these constitutive relations. Different microscopic
interactions can generate similar macroscopic
flows, while apparently small modifications
to the interaction rule can alter stability.
Hence, agreement with a measured velocity
field does not necessarily validate the
underlying closure approximation.

Robotic experiments are attractive because
they can provide both microscopic trajectories
and coarse-grained collective measurements.
In principle, we can characterize the
interaction rules, derive a continuum model,
and compare its predictions with independently
measured density and velocity fields.

Yet, this program requires sufficient
experimental scale. A continuum description
is most informative when a separation exists
between the microscopic interaction length
and the spatial scale of collective variation.
Many robotic collectives contain relatively
few agents, making this separation difficult
to establish.

\paragraph*{Experimental coverage.}

The extraction identifies 65 studies comparing
with kinetic or hydrodynamic theory.
Robotic micro- and nano-agents contribute 28,
vibration-driven robots 14,
and wheeled or legged robots eleven.
Five belong to the residual platform category,
three involve smarticles or granular robots,
two feedback-controlled colloids,
and one a modular system.
One study has no recorded platform class.

This distribution differs from that of
Vicsek-type models. Microscopic robotic
assemblies form the largest contributing
platform class, reflecting the importance
of fluid-mediated interactions and
collective transport at small scales.

The included microscopic experiments
investigate collective propulsion,
field-driven motion, hydrodynamic coupling,
and reconfigurable assemblies
\cite{Sun2021swarmingmicrodrople,Calabr2026interactiondrivena,Xie2020bioinspiredsoftmic,Pauer2021programmabledesign,Yigit2019programmablecollect,Das2023optoelectronictraje,McGovern2024multifunctionalchir,Sun2024multiplemagnetoopt,Liu2022bionicmorphological,Xu2023multiplecilialike,Yue2022wheellikemagnetic,Wang2018collectivebehavior,Chen2022programmablespatio,Xu2022collectivebehaviors,Baker2019shapeprogrammed3d,Goyal2025externallycontrolle,Aubret2018targetedassemblyan,Ahmed2021bioinspiredacousto,Yan2024programmingstructur,Yigit2020cohesiveselforgani,Hanson2025electrokineticpropu,Xu2023bioinspiredciliali,Ji2021ondemandassemblya,Xie2019programmablegenerat,Law2024navigationofmagnet,Yu2022adaptivepatternand,Ceron2026fluidictorqueenabl,Li2022vectorcontrolledwh}.

The extracted observables include particle
velocities, relative displacement,
angular motion, fluid-flow measurements,
cluster geometry, and changes in
collective organization. These quantities
allow comparison with theoretical descriptions
at several levels.

For example, a measured flow field can
be compared with a hydrodynamic calculation.
Particle trajectories can constrain
effective mobility or coupling.
Density and cluster statistics can
test predictions for collective
organization.

However, these comparisons are not
equivalent. Matching the response
to an imposed external field
tests a different theoretical
statement from predicting
spontaneous collective structure.

\paragraph*{Mechanical and macroscopic realizations.}

Vibration-driven robots provide a complementary
experimental regime in which continuum descriptions
may emerge from contact-mediated interactions
rather than fluid coupling
\cite{Scholz2018rotatingrobotsmove,Arbel2025amechanicalroutef,Baconnier2022selectiveandcollec,Casiulis2025ageometricconditio,Deblais2018boundariescontrolc,Boudet2022effectivetemperatur,Wang2025hyperuniformnetwork,Miyashita2007selfassembleofwat,Mayya2019nonuniformrobotde,Li2021programmingactivec,Zion2023morphologicalcomput,Wagner2026geometryinducedtra,Lvay2026collectivedynamics,Zion2021distributedonline}.

These studies include measurements
of spatial correlations,
mean-squared displacement,
velocity distributions,
collective oscillations,
and structural organization.

Their importance is that effective
continuum behavior may arise
from mechanically coupled
constituents whose microscopic
dynamics are not specified
by an alignment controller.
This provides an opportunity
to test continuum descriptions
against independently generated
collective motion.

Yet, dry mechanical systems
may possess frictional
and contact-mediated
interactions that do not
satisfy the assumptions
of conventional fluid
hydrodynamics. An appropriate
theoretical comparison
must retain the relevant
physical mechanisms.

Wheeled or legged robots offer
different access to continuum
behavior through measurements
of population density,
collective motion,
and spatial organization
\cite{Giardina2022collectivephototact,Fersula2026aggregatingswarmst,Valentini2015collectivedecision,Zhou2024anticipationpromote,Hamann2008spatialmacroscopic,Maffettone2024mixedrealityenviro,Aguilar2018collectiveclogcont,Kernbach2013adaptivecollective,Arvin2011imitationofhoneybe,Prasath2022dynamicsofcooperat,Fruchart2020phasetransitionsin}.

The main advantage is experimental control
of microscopic rules. The corresponding
limitation is that the same rules may
directly prescribe the collective
response being compared.

A particularly informative experiment
would derive a continuum prediction
from the implemented interaction
and test its spatial or temporal
dependence without separately
fitting every observable.

Additional implementations address
environment-mediated flows,
collective mechanical response,
and hydrodynamic interactions
\cite{Bo2026threephasesofodd,Newbolt2024flowinteractionsle,Boudet2025effectsofboundarie,Ren2025chlamydomonasinspir,Leech2026spincoupledhydrody,Heuthe2025tunablecolloidalsw,Guo2026afeedbackcontrolle,Luo2025flockingphasesepar,Savoie2019arobotmadeofrobo,Li2020programmingactiveg,Miyashita2008howmorphologyaffec,Fruchart2021nonreciprocalphase}.
Together, these studies demonstrate
that continuum and kinetic ideas
have been brought into contact
with several physically different
robotic systems.

\paragraph*{Agreement and experimental limitations.}

Of the 65 studies, 29 report
quantitative theory--experiment
agreement and 36 qualitative agreement.
None has an agreement classification
of ``none'' in the extracted field.

For D6, 24 studies are rated
low risk, 33 have some concerns,
and eight are high risk.
Thus, a substantial subset
provides relatively strong
quantitative comparisons.

Replication reporting remains
less complete. Twenty studies
are rated low risk in D1,
and 45 have some concerns.
Independent-run counts
are recorded in 24 studies.

Explicit code availability
is recorded for twelve studies
and data availability for sixteen.
One study was assessed from
an abstract alone. The remaining
64 were assessed on full text.

These indicators support
a cautiously positive assessment.
Kinetic and hydrodynamic theories
have been compared quantitatively
with several robotic platforms.
However, the evidence does not
establish that the same continuum
equations describe all these systems.

The distinction is particularly
important for hydrodynamic coupling.
A microscopic suspension interacts
through the surrounding fluid,
whereas a dry robotic collective
may interact through collisions,
sensing, or mechanically transmitted
forces. Similar density or
velocity fields need not imply
identical constitutive equations.

\paragraph*{What constitutes a decisive continuum test?}

A useful comparison should connect
microscopic parameters to
coarse-grained observables.

First, the interaction law
should be characterized
independently where possible.
For microscopic swimmers,
this may require measurements
of pairwise relative motion
or flow-mediated coupling.
For macroscopic robots,
the implemented sensing
and control rules can
provide the starting point.

Second, the continuum fields
should be measured over
a sufficient spatial range.
A density profile or
velocity field inferred
from a handful of robots
may be dominated by
finite-size effects.

Third, predictions should
be tested under changes
in the relevant control
parameters. The strongest
evidence would connect
one parameterization
to several observables:
for example, density
fluctuations, transport,
and response to confinement.

A continuum model that
reproduces one steady-state
density profile may be
useful. A model that
predicts the evolution
of density and polarization
after a perturbation,
without additional fitting,
provides a stronger test.

We therefore conclude
that kinetic and
hydrodynamic theories
form an important
quantitative component
of robotic active-matter
research. Their main
opportunity is the
connection between
programmable microscopic
interactions and
measured collective fields.

The present limitation
is the diversity of
physical mechanisms,
experimental scales,
and comparison protocols.
Progress will depend
on systematic measurements
of continuum observables
and clearer demonstrations
of where the corresponding
approximations remain valid.

\subsection{Active elasticity}
\label{sec:theory-active-elasticity}

Active elasticity addresses collective systems in which
internally generated forces modify mechanical response.
Unlike models focused primarily on particle trajectories,
elastic theories describe the relation between deformation,
stress, and the dynamics of coupled mechanical degrees
of freedom.

Active systems can exhibit nonreciprocal response
because their forces need not derive from an
equilibrium elastic energy
\cite{Scheibner2020odd}
\cite{Scheibner2020odd}.

The experimental challenge is to distinguish
an active constitutive property from
the response of one particular controller
or mechanical assembly.

For a linearized mechanical system,
the relation between generalized
displacements and forces may be written
\begin{equation}
F_a=\sum_b K_{ab}u_b,
\label{eq:theory-active-stiffness}
\end{equation}
where $K_{ab}$ is an effective
response matrix.

In a reciprocal passive system,
appropriate equilibrium assumptions
constrain the symmetry of this matrix.
Active mechanical coupling can
produce a nonreciprocal contribution.

The relevant experiment measures
the response coefficients and
their dependence on actuation,
geometry, or driving frequency.
Observing a moving mechanical
structure is not sufficient
to establish an active
elastic constitutive law.

\paragraph*{Experimental comparisons.}

The extraction identifies seven studies
comparing robotic experiments with
active elasticity.

Four use robotic metamaterials
or active solids, two belong
to the residual platform category,
and one uses vibration-driven robots.

The mechanically focused comparisons
include experiments on nonreciprocal
coupling, active elastic response,
and collective structural dynamics
\cite{Brandenbourger2021limitcyclesturnac,Binysh2026moreislessinunpe,Coulais2025activemetamaterials,Chen2020selfsensingmetamat}.

The extracted observables include
shear response,
odd elastic coefficients,
vibrational properties,
wave amplitude,
frequency,
and spatial decay.

These quantities permit
comparisons with predictions
that are more specific
than collective motion
alone.

For example, frequency-dependent
wave propagation can
constrain the dynamical
response of the structure.
A measured stiffness
or response tensor
can test predicted
nonreciprocal coupling.

Other experimental comparisons
involve mechanically driven
or programmed collective
systems
\cite{Mohanty2025multimodalmotionan,nr2022mindthegappredic,Karagzel2024selforganizedflock}.
Their extracted observables
include velocity,
collective order,
and spatial response.

However, these measurements
do not all test
active elasticity
at the same level.
A collective trajectory
may be compatible
with an active
elastic description
without independently
establishing its
constitutive coefficients.

We therefore distinguish
direct mechanical
response measurements
from more indirect
comparisons between
collective dynamics
and elastic models.

\paragraph*{Quantitative evidence.}

Two of the seven studies
report quantitative
theory--experiment agreement
\cite{Brandenbourger2021limitcyclesturnac,Chen2020selfsensingmetamat}.
The remaining five
report qualitative agreement.

For D6, two studies
are rated low risk,
three have some concerns,
and two are high risk.

Thus, despite the
physical specificity
of active elasticity,
the extracted evidence
contains only a
small quantitative
testing subset.

Replication documentation
is also limited.
One study receives
a low-risk D1 rating,
while six have
some concerns.
Independent-run counts
are reported in
two studies.

Explicit code availability
occurs in one study
and data availability
in two. One study
was assessed on
an abstract-only basis.

These limitations
do not diminish
the conceptual
importance of
active robotic
materials. They
indicate that the
available evidence
is not yet broad
enough to establish
general constitutive
behavior across
independent
experimental platforms.

\paragraph*{What should be tested next?}

A decisive active-elasticity
experiment should measure
mechanical response under
several independent
perturbations.

This includes the
possibility of
testing reciprocity
by comparing the
response to loading
at different
degrees of freedom.
Frequency-dependent
measurements can
then establish
whether an effective
elastic description
remains appropriate
over the relevant
dynamical range.

Boundary conditions
are particularly
important. A finite
active structure
may display
localized modes
or apparent
nonreciprocity
that depend on
its geometry.

The comparison
should therefore
distinguish
bulk constitutive
properties from
boundary-mediated
behavior.

The strongest
evidence would
use independently
characterized
actuation and
mechanical parameters
to predict
deformation fields,
response spectra,
and work over
cyclic loading.

Such measurements
could test whether
a continuum
active-elastic
description remains
valid when network
geometry or
system size changes.

We conclude that
robotic active
solids provide
particularly direct
access to the
mechanical predictions
of nonequilibrium
theory. Yet, the
quantitative evidence
remains concentrated
in a small number
of experiments.

The immediate
priority is therefore
independent reproduction
of measured response
coefficients and
controlled tests
across different
mechanical structures.

\subsection{Other statistical-mechanical frameworks}
\label{sec:theory-statistical}

The theoretical diversity of robotic active matter
extends beyond alignment, persistent-particle,
hydrodynamic, and elastic descriptions.
Some experiments use statistical-mechanical
concepts to characterize fluctuations,
configuration probabilities, transport,
collective response, or nonequilibrium
state selection.

The extraction groups these approaches
under ``other statistical mechanics''.
This is a residual theory category,
not a single mathematical framework.
Its contents should therefore be
interpreted through the actual
observables and mechanisms
reported by each study.

The category contains 25 studies.
Six involve wheeled or legged robots,
five vibration-driven systems,
four feedback-controlled colloids,
three microscopic robotic agents,
three smarticles or granular robots,
three residual platforms,
and one modular system.

Its broad distribution is noteworthy.
Statistical-mechanical descriptions
are being used in systems with
informational interactions,
mechanical contact,
field-mediated actuation,
and programmable feedback.

Yet, the presence of a common
theoretical classification does
not establish common assumptions
about conservation laws,
effective temperature,
detailed balance,
or interaction symmetry.

\paragraph*{Configuration statistics and collective dynamics.}

Smarticle and granular-robot
experiments include comparisons
involving collective transport,
configuration probabilities,
and fluctuation-related
quantities
\cite{Li2018supersmarticlealo,Chvykov2020lowrattlingapred,Li2020programmingactiveg}.

These systems are useful
because internal actuation
and mechanical constraints
jointly determine the
accessible collective
configurations.

A statistical description
can therefore connect
the observed distribution
of configurations to
the underlying dynamics.
However, one must not
assume that a stationary
distribution has the
same interpretation
as an equilibrium
Boltzmann distribution.

Active driving can
maintain probability
currents even when
the measured distribution
appears stationary.
Consequently, the
configuration probability
alone may not determine
the nonequilibrium
dynamics.

Feedback-controlled
colloids provide
complementary tests
through interaction
potentials, rotational
response, and
configuration statistics
\cite{Yan2014colloidalsuperstruc,Buerle2020formationofstable,Knippenberg2022roleofcohesionin,Lffler2021behaviordependentc}.

These studies include
quantities such as
pair potentials,
rotational order,
susceptibility,
relaxation time,
and flow statistics.

The availability
of feedback control
permits effective
interactions to
be changed directly.
This offers an
opportunity to
compare statistical
responses under
different microscopic
rules.

Yet, a feedback-mediated
interaction may depend
on the measured state
and the control
implementation.
A conventional
equilibrium potential
is therefore not
automatically an
adequate description.

\paragraph*{Other robotic realizations.}

Microscopic robotic systems
include comparisons involving
collective assembly,
fluctuations, and
effective stochastic
descriptions
\cite{Yang2022emergentmicroroboti,Aubret2018targetedassemblyan,Raghavendra2026programmablepersist}.

Vibration-driven experiments
provide additional access
to collective statistics,
fluctuation spectra,
and configuration-dependent
dynamics
\cite{Baconnier2022selectiveandcollec,Boudet2025nonmarkoviandynami,Mayya2019nonuniformrobotde,Zhang2024proteinfoldinginsp,Li2021programmingactivec}.

These experiments illustrate
an important strength
of mechanical active
systems: statistical
behavior is not
necessarily prescribed
by an explicit
interaction algorithm.

The corresponding
theoretical comparison
can reveal whether
simple stochastic
descriptions capture
the observed collective
statistics.

Wheeled or legged
systems offer
a different approach
through explicitly
programmed interactions
and measured collective
order
\cite{Sankar2026trendformationwith,Turgut2008selforganizedflocka,Nishinari2006modellingofselfdr,Spears2004distributedphysics,Barone2024experimentalandnum,Rezeck2021flockingsegregative}.

Other specialized
platforms investigate
chiral collective
response, active
ordering, and
nonequilibrium
transport
\cite{Li2024memoryinducedspont,Bo2026threephasesofodd,Das2025roleofkinematicco,Kaeser2025individualandcolle}.

These implementations
provide useful
examples of
statistical-mechanical
reasoning in
robotic experiments.
However, the
specific theoretical
assumptions differ,
so the studies
cannot be combined
into one test
of a common model.

\paragraph*{Evidence and limitations.}

Sixteen of the 25 studies
report quantitative
agreement, while nine
report qualitative
agreement.

The D6 ratings are
low risk for thirteen,
some concerns for eleven,
and high risk for one.
Thus, this category
contains a comparatively
high proportion of
quantitative and
lower-risk comparisons.

Replication evidence
is less complete.
Seven studies receive
low-risk D1 ratings,
seventeen have
some concerns,
and one is high risk.
Independent-run counts
are recorded in
eight studies.

Explicit code availability
occurs in four studies
and data availability
in three.
Two studies were
assessed from
abstracts alone.

The main conclusion
is that statistical-mechanical
methods can describe
several robotic
collective observables
quantitatively.
However, the residual
category does not
identify one general
theory whose validity
can be assessed
across all
included systems.

The strongest
comparisons are
those that specify
a particular
statistical model,
measure its parameters,
and test the resulting
predictions across
independent conditions.

For nonequilibrium
systems, it is
especially important
to examine dynamical
observables as well
as stationary
distributions.
Two different
processes can
produce similar
configuration
statistics while
exhibiting different
probability currents
and temporal
correlations.

Hence, future
experiments should
combine configuration
measurements with
fluctuation,
relaxation,
and response
statistics.
Such measurements
would help identify
where an effective
statistical description
is predictive,
rather than merely
a compact fit
to observed data.

\subsection{Agent-based simulation as a comparator}
\label{sec:theory-agent-based}

Agent-based simulation is the most frequently
recorded theoretical comparator in the evidence base.
Its popularity is unsurprising. A simulation can
represent the individual rules of a robotic system,
including sensing, communication, actuation,
and interactions. It can also explore conditions
that are difficult or expensive to realize
in hardware.

Yet, an agent-based simulation occupies a
different scientific position from a
coarse-grained physical theory.
It may reproduce the implemented controller
almost directly. Agreement with experiment
then demonstrates that the simulation
captures relevant aspects of the hardware,
but it does not necessarily identify
a universal collective mechanism.

This distinction is particularly important
for a review of active-matter physics.
The value of an experimental comparison
depends on what the model predicts
beyond the rules already implemented.

The extraction identifies 275 studies
using agent-based simulation.
Wheeled or legged robots contribute 127,
other platforms 52,
vibration-driven robots 32,
and robotic micro- and nano-agents 30.
Feedback-controlled colloids contribute ten,
modular robots nine,
robotic active solids eight,
and smarticles or granular robots six.
One study has no platform classification.

Thus, agent-based simulation is represented
across nearly the entire robotic taxonomy.
The breadth is an important methodological
strength. It also makes the category
highly heterogeneous.

\paragraph*{What simulations can establish.}\hfill\break % layout only
At their most useful, agent-based simulations
provide a bridge between microscopic
implementation and collective observables.
A model can incorporate the interaction
rules and predict group-level behavior
under changes in agent number,
noise, geometry, or driving.

The robotic literature includes
such comparisons in flocking,
aggregation, transport,
and programmable collective
organization
\cite{Turgut2008modelingphasetrans,Navarro2008anapproachtoflock,Garnier2005aggregationbehaviou,Guzzi2014bioinspiredobstacle,Cazenille2016automatedoptimisati,Vsrhelyi2018optimizedflockingo,Ferrante2012selforganizedflock,Mijalkov2016engineeringsensoria,Berlinger2021selforganizedevasi,Fruchart2022nonreciprocalphase,Zhou2024anticipationpromote,Barone2024experimentalandnum,Paramanick2025spontaneousemergenc,Zhang2026emergentnonhermiti}.

Other comparisons investigate
collective motion,
configuration,
and interactions
under programmable control
\cite{Yamada2011circularmotiongene,Yang2023collectivechaseem,Tran2022frontierledswarmin,Giardina2022collectivephototact,Vega2023simulatelessexpec,Vega2023indirectswarmcontr,Vega2024agentbasedemulatio,Gao2023robustdistributedc,Lee2022distributedcohesive,Baumann2024hybridflockforma,Wu2024adversarialimitatio,Varadharajan2026programmableflocks,Pan2026exploringthephysic,Sankar2026trendformationwith,Verdoucq2025flockingphasetrans,Jond2026positionbasedflock}.

These systems allow the
experimenter to ask
whether simulated
collective behavior
survives the transition
to physical hardware.

This is a nontrivial
question. Real robots
possess actuator
limitations, sensing
errors, latency,
and physical contacts
that may be simplified
or omitted in simulation.

Consequently,
agreement with hardware
can validate aspects
of the simulation
as an engineering
model. It can
also reveal which
physical constraints
must be included
to reproduce
collective dynamics.

However, reproducing
the trajectory of
a programmed swarm
is a weaker test
of active-matter theory
than predicting
collective fluctuations
or transitions
from independently
characterized rules.

\paragraph*{Simulations of physical interactions.}

For mechanically driven
and microscopic systems,
agent-based simulation
can represent physical
interactions that
are not directly
programmed by
a controller.

The included studies
cover collision-mediated
dynamics,
mechanical aggregation,
fluctuations,
and collective transport
\cite{Welch2018aggregatesuppressio,Welch2019activereorientation,Xi2024emergentintelligenc,GuzmnLastra2026emergenceofbeating,Edwards2026robusttopologically,Musacchio2026activeembracemente,Arbel2025amechanicalroutef,Baconnier2022selectiveandcollec,Casiulis2025ageometricconditio,Deblais2018boundariescontrolc,Boudet2022effectivetemperatur,Wang2025hyperuniformnetwork,Musacchio2026fluidizationinduced,Zion2023morphologicalcomput,Wagner2026geometryinducedtra,Zion2021distributedonline}.

Microscopic robotic experiments
provide further comparisons
with models of field-driven
and interaction-mediated
collective dynamics
\cite{Sun2021swarmingmicrodrople,Calabr2026interactiondrivena,Xie2020bioinspiredsoftmic,Pauer2021programmabledesign,Chen2024synchronousandfull,Yigit2019programmablecollect,McGovern2024multifunctionalchir,Feng2023weakionexchangeba,Sun2024multiplemagnetoopt,Liu2022bionicmorphological,Wang2018collectivebehavior,Baker2019shapeprogrammed3d,Goyal2025externallycontrolle,Aubret2018targetedassemblyan,Yan2024programmingstructur,Yigit2020cohesiveselforgani,Hanson2025electrokineticpropu,Law2024navigationofmagnet,Ceron2026fluidictorqueenabl,Li2022vectorcontrolledwh}.

In these systems,
the microscopic
interaction model
may itself be
a scientific
hypothesis.
A simulation
based on measured
forces or
hydrodynamic coupling
can therefore
provide a substantive
test of physical
mechanisms.

Still, the strength
of the conclusion
depends on parameter
identification.
If interaction
parameters are
adjusted to reproduce
every experimental
configuration,
the model is
primarily descriptive.

A stronger comparison
uses independently
measured parameters
to predict how
the collective
responds to changes
in density,
confinement,
or driving.

\paragraph*{Mechanical assemblies and active materials.}

Agent-based simulations
also play an important
role in systems
with changing connectivity
or mechanically
coupled components.

The included literature
contains modular robots,
smarticles,
and active-solid
experiments
\cite{Vardhan2022glidersinshapecha,Savoie2019arobotmadeofrobo,Li2020programmingactiveg,Hu2026diatominspired1di,Rogowski2026coordinatedselfass,Kaeser2025individualandcolle,Bhattacharjee2026collectiveassembly,Li2019particleroboticsba,Baconnier2022tensioncontrolledc,Rezanejad2025programmableentangl,Brandenbourger2019nonreciprocalrobot,Brandenbourger2021limitcyclesturnac,Binysh2026moreislessinunpe,Coulais2025activemetamaterials,Chen2020selfsensingmetamat}.

The relevant observables
include displacement,
assembly geometry,
mechanical response,
configuration statistics,
and collective
deformation.

These quantities
can establish whether
a discrete model
captures the behavior
of a coupled structure.
However, agreement
with a detailed
simulation does not
automatically validate
a proposed continuum
constitutive law.

The distinction is
particularly important
for active elasticity.
A discrete mechanical
model may reproduce
the response of
a finite network,
while a continuum
description requires
additional assumptions
about scale separation
and effective
material properties.

Thus, agent-based
simulation is often
a necessary intermediate
step. It should not
be confused with
the final physical
explanation.

\paragraph*{Quantitative evidence and reproducibility.}

Among the 275 studies,
96 report quantitative
agreement and 171
qualitative agreement.
Seven are classified
as having no agreement
comparison, and one
has a missing
agreement value.

The D6 ratings are
low risk for 62 studies,
some concerns for 110,
and high risk for 103.

This distinction is
important. Although
agent-based simulation
is the most frequent
comparison category,
fewer than one quarter
of its studies receive
low-risk D6 ratings.

Replication information
is also incomplete.
Seventy-eight studies
are rated low risk
in D1, 190 have
some concerns,
and seven are high risk.

Independent-run counts
are recorded for
107 studies.
Explicit code availability
occurs in 45,
and data availability
in 37.
Eight studies were
assessed from
abstracts alone.

The availability
of simulation code
is particularly
important for this
framework.
A published model
description may omit
implementation details
that materially affect
collective behavior,
including update order,
numerical integration,
collision handling,
and stochastic
sampling.

For robotic systems,
hardware-controller
code may be equally
important.
A discrepancy between
simulation and experiment
can result from
an omitted physical
interaction or
a difference in
the implemented
control rule.

Consequently,
reproducible
theory--experiment
comparison should
document both
the computational
model and the
experimental
implementation.

\paragraph*{From simulation agreement to physical explanation.}

We distinguish three scientifically
different uses of agent-based simulation.

A hardware-oriented
simulation reproduces
the robot's control
and actuation.
Its main purpose
is to test whether
the implementation
behaves as expected.

A phenomenological
simulation reproduces
selected collective
patterns.
It provides a
compact description,
but different
microscopic models
may generate similar
collective behavior.

A mechanistic
simulation predicts
collective observables
from independently
characterized
interactions.
This is the
strongest basis
for testing a
physical explanation.

The extracted
theory-comparison
labels do not
systematically
separate these
three uses.
Hence, the 96
quantitative agreements
cannot be interpreted
as 96 independent
mechanistic validations.

This limitation
also affects
comparisons among
theoretical frameworks.
An agent-based
simulation may
implement an
active Brownian
particle model,
a Vicsek-type
rule, or a
discrete elastic
network.
Because the
framework field
is multi-label,
these comparisons
may be represented
in several categories.

The appropriate
interpretation
therefore concerns
the mathematical
model and the
tested observable,
not the presence
of a simulation
label alone.

We conclude that
agent-based simulation
is indispensable
for connecting
robotic implementation
to collective dynamics.
Its scientific
strength is the
ability to vary
microscopic rules
and compare their
consequences with
physical experiments.

Yet, its prevalence
should not be
mistaken for
independent theoretical
validation.
The most informative
studies use simulation
to make predictions
that could fail,
rather than merely
to reproduce
the behavior
of a controller.

\subsection{What constitutes a decisive theoretical test?}
\label{sec:theory-decisive}

The comparison across theoretical frameworks
reveals a consistent limitation.
Robotic active matter is highly programmable,
but quantitative predictivity is less
widespread than implementation and
qualitative agreement.

This does not mean that qualitative
experiments lack value.
A new robotic platform can establish
that an interaction or collective
state is physically realizable.
It may also reveal unexpected
behavior that motivates
new theory.

Yet, the requirements for
testing an established
physical model are different.
We must determine which
observations the model
explains, which parameters
were adjusted, and which
predictions were evaluated
independently.

\paragraph*{Fitted parameters versus predicted observables.}

Suppose a theoretical model predicts
an observable $O$ as a function of
control parameters $\boldsymbol{\lambda}$
and model parameters $\boldsymbol{\theta}$:
\begin{equation}
O_{\mathrm{th}}
=
F(\boldsymbol{\lambda},
  \boldsymbol{\theta}).
\label{eq:theory-prediction}
\end{equation}

A comparison becomes more informative
when $\boldsymbol{\theta}$ is
constrained independently,
rather than adjusted
separately for every
experimental condition.

If the parameters are
fitted to reproduce
one measured curve,
agreement with that
curve establishes
descriptive adequacy.
It does not
independently validate
the model's prediction.

By contrast,
a model fitted
under one condition
and evaluated under
another can test
transferability.
An even stronger
experiment determines
the microscopic
parameters independently
and predicts the
collective response
without fitting
to that response.

This distinction
is relevant to
all the frameworks
considered here.

For Vicsek-type
models, the
implemented alignment
rule can be known,
while collective
fluctuations remain
to be predicted.

For active Brownian
particles, isolated-agent
motility statistics
can constrain
the parameters
before density-dependent
collective behavior
is measured.

For hydrodynamic
theories, microscopic
interactions can
inform constitutive
coefficients before
collective fields
are compared.

For active elasticity,
measured response
coefficients can
be used to predict
deformation or
wave propagation
under a different
loading condition.

These approaches
would turn
programmability
into independent
experimental constraint.

\paragraph*{Parameter sweeps and competing predictions.}

A single successful comparison
rarely distinguishes
between alternative
models. Several
theoretical descriptions
can reproduce one
collective trajectory,
density profile,
or steady-state
configuration.

A stronger experiment
varies an independently
controlled parameter
and compares the
resulting response
with competing
predictions.

For example,
changing particle
density can test
clustering models,
while varying
interaction range
can test alignment
theories.
Changing persistence
can distinguish
different stochastic
motility descriptions.
Controlled loading
can test constitutive
models of active
solids.

The most informative
parameter changes
are those for
which different
models predict
different outcomes.

A parameter sweep
is therefore not
merely a way to
produce more data.
It is a method
for making
theoretical
assumptions
experimentally
discriminable.

This requires
careful calibration.
Changing robot
speed may also
alter collision
rates, motor
noise, and
turning dynamics.
Similarly, increasing
a magnetic field
can modify both
propulsion and
interparticle forces.

Where control
parameters are
physically coupled,
the model must
account for
that coupling.
Otherwise,
a quantitative
match may depend
on compensating
errors.

\paragraph*{Several observables from one parameterization.}

Another demanding
test is simultaneous
agreement with
different collective
observables.

A model may reproduce
mean-squared displacement
without reproducing
orientation correlations.
It may predict
a steady-state
density profile
while failing
to describe
relaxation after
a perturbation.

Likewise,
a mechanical model
may reproduce
a deformation
pattern without
predicting the
measured frequency
response.

These discrepancies
are scientifically
informative.
They identify
which approximations
are adequate
and which
physical mechanisms
are missing.

Consequently,
we should not
treat a mismatch
as a failed
experiment.
A well-characterized
failure of a
theoretical prediction
can provide
stronger evidence
than a visually
convincing agreement
obtained through
extensive fitting.

For robotic
active matter,
the most useful
combinations often
include both
microscopic and
collective observables.

For example,
isolated-agent
velocity correlations
can be compared
with collective
transport,
while pairwise
interaction statistics
can be related
to clustering
or alignment.

This provides
a direct connection
between mechanisms
and emergent
behavior.

\paragraph*{Experimental scale and finite-size effects.}

Many theoretical
predictions concern
collective states
that become sharply
defined only
at large system
size.

Robotic experiments
are finite,
sometimes containing
only a handful
of agents.
This need not
invalidate their
scientific contribution.
Small systems
can provide
precise measurements
of interactions
and local dynamics.

However,
they cannot
independently
establish
thermodynamic-limit
behavior without
additional analysis.

For example,
persistent polarization
in a small robot
collective does
not determine
the scaling
of orientational
fluctuations.
A cluster that
forms in a
confined arena
does not
necessarily
represent
phase separation.

Similarly,
a localized
mechanical mode
in a short
active network
does not
automatically
establish a
bulk constitutive
property.

Where theories
predict scaling,
we need
measurements
at several
system sizes.
The corresponding
boundary conditions
must also be
characterized.

This connects
directly to
the review's
D3 finite-size
and boundary
risk-of-bias
domain.
A theoretical
comparison can
be numerically
accurate while
remaining limited
to one geometry
or system size.

The correct
scientific conclusion
should preserve
that domain
of applicability.

\paragraph*{Independent replication and research groups.}

A general physical
claim becomes
stronger when
the relevant
prediction is
reproduced independently.

This independence
can concern
the experimental
hardware,
the research group,
the analysis
method, or
the physical
implementation
of the interaction.

For robotic
active matter,
cross-platform
replication is
especially valuable.
A theory that
predicts comparable
collective behavior
in mechanically
driven particles,
programmable
mobile robots,
and microscopic
assemblies may
have wider
applicability
than one tested
on a single
hardware design.

Yet, such
comparison requires
careful identification
of the shared
physical assumptions.
Different platforms
need not use
identical components,
but their relevant
dimensionless
parameters and
observables
must be comparable.

The extracted
dataset records
publication metadata,
platform classes,
and theoretical
comparisons.
It does not
provide a normalized
mapping of independent
experimental groups.
We therefore
avoid treating
different publications
as automatically
independent
replications.

Likewise,
multiple reports
from closely
related hardware
systems need
not constitute
independent
experimental
confirmation.

A future living-review
update could
record experimental
lineage,
hardware reuse,
and independent
replication explicitly.
This would improve
certainty assessments
beyond publication
counts alone.

\paragraph*{Reporting requirements for reproducibility.}

The review's quality
indicators identify
several recurring
limitations across
theoretical frameworks.

Independent-run
counts are often
missing.
Code and data
availability are
incomplete.
Finite-size
and boundary
effects are
not always
characterized.

These omissions
matter because
theory--experiment
agreement is
typically evaluated
through estimated
observables.
Without sufficient
replication,
we cannot
independently
assess their
statistical
uncertainty.

Without accessible
trajectories or
measurement data,
it can be difficult
to verify order
parameters,
correlation functions,
or fitted
scaling exponents.

Without controller
and simulation
implementations,
the comparison
may depend on
unreported details
of the interaction
rule or numerical
procedure.

These requirements
are not identical
across platforms.
A mechanical
active solid
requires appropriate
response calibration,
while a feedback-controlled
colloidal system
requires an accurate
description of
the measurement
and feedback
process.

Nevertheless,
the underlying
principle is shared:
the experimental
conditions necessary
to reproduce
the prediction
must be documented.

\paragraph*{From agreement to explanatory power.}

We can now state
the central result
of the theoretical
synthesis.

Agent-based simulation
is the dominant
comparison method
in the extracted
literature.
Its breadth
reflects its
flexibility,
but many comparisons
are qualitative
or closely tied
to the implemented
robotic controller.

Vicsek-type
models provide
directly programmable
alignment rules.
Their experimental
realizations are
numerous relative
to their low-risk
quantitative tests.

Active Brownian
particle descriptions
have a comparatively
strong quantitative
presence, particularly
in mechanically driven
systems.
Their advantage
is the connection
to measurable
single-agent
motility statistics.

Run-and-tumble
descriptions remain
represented by
only a few
experimental comparisons.
They offer
a useful opportunity
for distinguishing
different stochastic
motility mechanisms.

Kinetic and
hydrodynamic theories
connect microscopic
interactions
to collective
fields.
Their quantitative
evidence is
substantial but
distributed across
physically different
platforms and
interaction mechanisms.

Active elasticity
offers specific
mechanical predictions
that are directly
accessible to
robotic experiments.
However, independently
replicated quantitative
tests remain limited.

Other statistical-mechanical
frameworks have
provided several
quantitative comparisons,
but their residual
classification
does not represent
one unified
theoretical model.

Across all these
frameworks,
the strongest
scientific opportunity
is the same:
programmability
can be used
to construct
experiments
that discriminate
between physical
explanations.

The goal should
not be to reproduce
the greatest
number of
collective patterns.
It should be
to determine
which predictions
survive controlled
changes in
interaction,
activity,
geometry,
and scale.

We therefore
regard a
decisive robotic
active-matter
experiment as
one in which
the microscopic
implementation
is characterized,
the collective
observables
are measured
quantitatively,
and the theoretical
prediction is
evaluated under
conditions not
used to determine
its parameters.

Independent
reproduction
then strengthens
the claim
that the
result reflects
a general
physical mechanism
rather than
one particular
robotic implementation.

This standard
also identifies
where robotics
can contribute
most to
active-matter
physics.
A robot is
not merely
a visible
realization
of a
theoretical
particle.
It can be
an instrument
for changing
the assumptions
of a model
one at
a time.

Used in
this way,
robotic
collectives
allow us
to move
from demonstrating
nonequilibrium
phenomena
toward testing
the theories
that explain
them.
